\documentclass[10pt,a4paper]{amsart}
\usepackage[margin=19mm]{geometry}
\usepackage{amsmath,amssymb,mathtools}
\usepackage[scr=rsfs,cal=euler]{mathalfa}
\usepackage{xfrac}
\usepackage[unicode,hidelinks,bookmarks=false]{hyperref}
\newcommand{\ie}{i.e.\ }
\newcommand{\cf}{cf.\ }
\newcommand{\resp}{respectively\ }
\newcommand{\Subsection}{\subsection}
\providecommand{\nobreakdash}{\nobreak\hbox{-}}
\setlength{\emergencystretch}{2em}
\multlinegap0pt
\usepackage{amssymb}
\usepackage{bm}
\usepackage{tikz}
\usepackage{xcolor}
\usepackage{mathtools}
\usepackage{float}
\newtheorem{remark}{Remark}
\newtheorem{proposition}{Proposition}
\newcommand{\Union}{\bigcup}	
\newcommand{\hyper}[1]{{\bf #1}}
\newcommand{\inc}{\subseteq}
\newcommand{\inter}{\cap}
\newcommand{\restrconstr}[2]{{#1}_{{}^\lceil {#2}}}
\newcommand{\set}[1]{\{#1\}}
\newcommand{\setc}[2]{\set{#1 \mid #2}}
\title{Tridendriform algebras on hypergraph polytopes, the other way around}
\author{Pierre-Louis Curien, B\'er\'enice Delcroix-Oger, Jovana Obradovi\'c}
\date{Source: arXiv:2606.17755v2 (CC BY 4.0)}
\begin{document}
\maketitle

\noindent\emph{Source and licence.} Tridendriform algebras on hypergraph polytopes, the other way around. Version arXiv:2606.17755v2 (CC BY 4.0), distributed by arXiv under the Creative Commons Attribution 4.0 International licence, http://creativecommons.org/licenses/by/4.0/, as recorded in the arXiv licence metadata for this exact version and shown on the abstract page https://arxiv.org/abs/2606.17755v2. Full licence notice: \texttt{licenses/arxiv-2606.17755-CC-BY-4.0.txt}.\par\smallskip

\noindent\textit{Editorial scope.} Counted unit: the article's proposition shuffle-restriction (source0.tex lines 478--486). The article's complete original proof of it is delivered with it (source0.tex lines 487--536, one \texttt{proof} environment of 50 lines). No mathematics is written, repaired or re-derived: the statement and the proof are the article's own lines, in the article's own notation, and no retained line is substituted or re-flowed.  Retained uncounted as the source-local context the proof needs: lines 298--300.  Nothing was omitted from any retained span, so the package carries no omission record.  No retained line contains a citation, so the wrapper carries no bibliography and no bibliography entry.  No retained line contains a figure environment, an image-inclusion command or drawing code, so no image is delivered and the package declares no asset dependency.  Editorial formatting only, in addition: an AMS \texttt{amsart} preamble that restates the article's own declarations and macros used by the retained lines, namely the environments \texttt{remark}, \texttt{proposition} on separate counters, together with the standard packages those lines need.  The retained environments print in this document as Remark 1, Proposition 1 in order of appearance, whereas the article prints the same items with the numbering of its own full document.  The complete native source of the article as distributed in the arXiv e-print for this exact version is bundled unchanged as \texttt{sources/203031/source0.tex}.
\begin{remark}\label{rest1}
For connected subsets $\emptyset \neq L\subseteq K$ and $\emptyset\neq K\subseteq H$  and a construct $C:{\bf H}$,  it holds that ${(C_{\lceil_K})}_{\lceil_L}=C_{\lceil_L}$.
\end{remark}

\begin{proposition} \label{shuffle-restriction}
Let $\delta = (\{C_a:{\bf H}_a\,|\, a\in A\},{\bf H})$  be a delegation and let $\emptyset\neq K\subseteq H$ be a connected subset of $H$. For a subset $B\inc A$, set  $B_1:=\setc{b\in B}{X_b\inter K\neq\emptyset}$ and
$B_2:=\setc{b\in B}{X_b\inter K=\emptyset}$. Then the following properties hold.
\begin{itemize}
\item[{\em (P$_1$)}] {\it If $B_1\neq\emptyset$, then $\restrconstr{(\ast_B(\delta))}{K}=\restrconstr{(\ast_{B_1}(\delta))}{K}$.}
\item[{\em (P$_2$)}] {\it If $B_1=\emptyset$, then $\restrconstr{(\ast_B(\delta))}{K}=\ast(\restrconstr{\delta}{K})$.}
\item[{\em (P)}] The shuffle product is compatible with restriction, i.e., $\restrconstr{(\ast(\delta))}{K}=\ast(\restrconstr{\delta}{K})$.
\end{itemize}
\end{proposition}

\begin{proof} We shall prove   the three  properties   together, by induction on the size of ${\bf H}$.  

\smallskip

 We   prove  (P$_1$) first. We set $X_B=\Union_{b\in B}X_b$ and $X_{B_1}=\Union_{b\in B_1}X_b$, inducing
$$\begin{array}{rcl}

\hyper{H},X_B & \leadsto & \hyper{H}_1,\ldots,\hyper{H}_n\\ 
\hyper{H},X_{B_1} & \leadsto & \hyper{H}^1_1,\ldots,\hyper{H}^1_m \\
\hyper{H}_K,X_{B}\inter K & \leadsto &\hyper{H}_1^K,\ldots,\hyper{H}_p^K,
\end{array}$$
as well as (a unique) $\varphi:\set{1,\ldots,p}\rightarrow\set{1,\ldots,n}$, such that $H_k^K\inc H_{\varphi(k)}$ for all $1\leq k\leq p$. We observe that, by construction, $X_B\cap K=X_{B_1}\cap K$. Therefore, the third decomposition also induces (a unique) $\psi:\set{1,\ldots,p}\rightarrow\set{1,\ldots,m}$, such that $H_k^K\inc H^1_{\psi(k)}$ for all $1\leq k\leq p$. The left-hand side of ($P_1$) unfolds as follows:  
 $$\begin{array}{rcl}
\restrconstr{(\ast_B(\delta))}{K} & =&(X_ B(\ast(\delta\lceil_{H_1}),\dots,\ast(\delta\lceil_{H_n})))\lceil_{K} \\[0.1cm]
&=&(X_B\cap K)(\ast(\delta\lceil_{H_{\varphi(1)}})\lceil_{H^K_1}, \dots ,\ast(\delta\lceil_{H_{\varphi(p)}})\lceil_{H^K_p}) \\[0.1cm]
&=& (X_B\cap K) (\ast((\delta\lceil_{H_{\varphi(1)}})\lceil_{H^K_1}), \dots ,\ast((\delta\lceil_{H_{\varphi(p)}})\lceil_{H^K_p})) \\[0.1cm]
&=& (X_B\cap K)(\ast(\delta\lceil_{{H^K_1}}), \dots ,\ast(\delta\lceil_{{H^K_p}})), 
\end{array}$$
where we have successively used the definition of shuffle product, the definition of restriction, the induction hypothesis, and Remark \ref{rest1}. On the other hand, the right-hand side is computed in an analogous way, replacing $B$ by $B_1$, and $\varphi$ by $\psi$. Since $X_B\cap K=X_{B_1}\cap K$, we thus obtain $$\restrconstr{(\ast_B(\delta))}{K} =(X_{B_1}\cap K)(\ast(\delta\lceil_{{H^K_1}}), \dots ,\ast(\delta\lceil_{{H^K_p}}))=\restrconstr{(\ast_{B_1}(\delta))}{K}.$$
 
\smallskip
 
For the proof of  the property (P$_2$) , we consider only the decomposition $\hyper{H},X_B   \leadsto   \hyper{H}_1,\ldots,\hyper{H}_n$. We notice that the assumption $B_1=\emptyset$ amounts to $X_b\cap K=\emptyset$, which entails that there exists a unique $1\leq i\leq n$ such that $K\subseteq H_i$. The left-hand side now calculates as 
 $$\begin{array}{rcl}
\restrconstr{(\ast_B(\delta))}{K} & =&(X_ B(\ast(\delta\lceil_{H_1}),\dots,\ast(\delta\lceil_{H_n})))\lceil_{K} \\[0.1cm]
&=& (\ast (\delta\lceil_{H_i}))\lceil_K \\[0.1cm]
&=& \ast ((\delta\lceil_{H_i}) \lceil_K) \\[0.1cm]
&=& \ast (\delta\lceil_{K}),
\end{array}$$
where we followed exactly the same steps as above.  

\smallskip
 
Finally,
using the properties  (P$_1$) and  (P$_2$), we calculate
$$\begin{array}{l} \restrconstr{(\ast(\delta))}{K} \\
 =  \Sigma_{\emptyset\neq B\inc A} q^{|B|-1} \restrconstr{(\ast_B(\delta))}{K}\\
 =  \Sigma_{\emptyset\neq B_1\inc A_1,B_2\inc A_2} q^{|B_1|+|B_2|-1} \restrconstr{(\ast_B(\delta))}{K} +
\Sigma_{\emptyset\neq B_2\inc A_2} q^{|B_2|-1} \restrconstr{(\ast_B(\delta))}{K}\\
 =  \Sigma_{\emptyset\neq B_1\inc A_1} q^{|B_1|-1}(\Sigma_{B_2\inc A_2} q^{|B_2|})\restrconstr{(\ast_{B_1}(\delta))}{K} +
(\Sigma_{\emptyset\neq B_2\inc A_2} q^{|B_2|-1}) \ast(\restrconstr{\delta}{K}).
\end{array}$$
Now we observe (using $q=-1$) that 
$$\begin{array}{l}
\Sigma_{B_2\inc A_2} q^{|B_2|} = (1+q)^{|A_2|}=0\\
\Sigma_{\emptyset\neq B_2\inc A_2} q^{|B_2|-1} = \frac{(1+q)^{|A_2|}-1}{q}= 1,
\end{array}$$
and we conclude that $ \restrconstr{(\ast(\delta))}{K} = \Sigma_{\emptyset\neq B_1\inc A_1}0 \:\restrconstr{(\ast_{B_1}(\delta))}{K} +
\ast(\restrconstr{\delta}{K}) = \ast(\restrconstr{\delta}{K})$.
\end{proof}

\end{document}
